\section{Positive-diagonal Matignon stability and orbit reduction}
\label{sec:framework}

For the commensurate Caputo system
\[
{}^C D_t^\alpha x=Ax,
\qquad 0<\alpha<1,
\]
asymptotic stability is equivalent to
\[
\sigma(A)\subset
\Sigma_\alpha
:=
\{z\neq0:|\arg z|>\alpha\pi/2\}
\]
\citep{Matignon1996,BrandiburGarrappaKaslik2021}. Define
\[
\mathcal F_\alpha^{(n)}
=
\{A:\sigma(DA)\subset\Sigma_\alpha
\ \forall D\succ0\text{ diagonal}\},
\]
\[
\mathcal D_H^{(n)}
=
\{A:DA\text{ Hurwitz }\forall D\succ0\},
\qquad
\mathcal P_\alpha^{(n)}
=
\mathcal F_\alpha^{(n)}\setminus\mathcal D_H^{(n)}.
\]
All interiors are taken in the Euclidean topology of \(\mathbb R^{n\times n}\). Because
\[
DA=D(AD)D^{-1},
\]
left and right positive-diagonal multiplication are spectrally equivalent. Our Hurwitz sign convention corresponds to Cain's positive-stability convention after replacing his matrix \(M\) by \(-A\) \citep{Cain1976}.

A real matrix is a P-matrix if every principal minor is positive and a \(P_0\)-matrix if every principal minor is nonnegative. Assume \(-A\) is strict P and, in dimension three, write
\[
p_i=-a_{ii}>0,
\qquad
m_{ij}=\det A[\{i,j\}]>0,
\qquad
q=-\det A>0.
\]
Define the positive-left-diagonal invariants
\[
\beta_{12}=\frac{m_{12}}{p_1p_2},
\quad
\beta_{13}=\frac{m_{13}}{p_1p_3},
\quad
\beta_{23}=\frac{m_{23}}{p_2p_3},
\quad
\kappa=\frac{q}{p_1p_2p_3}.
\]

\begin{lemma}[Positive-diagonal orbit reduction]
\label{lem:simplex-reduction}
Let
\[
D=\operatorname{diag}(d_1,d_2,d_3)\succ0,
\qquad
s=\sum_{i=1}^3p_id_i,
\qquad
x_i=\frac{p_id_i}{s}.
\]
Then \(x\in\Delta_2^\circ=\{x_i>0:\sum_i x_i=1\}\), and
\[
\det(\lambda I-DA)
=
\lambda^3+s\lambda^2+s^2B_\beta(x)\lambda
+s^3\kappa x_1x_2x_3,
\]
where
\[
B_\beta(x)
=
\beta_{12}x_1x_2+
\beta_{13}x_1x_3+
\beta_{23}x_2x_3.
\]
Conversely every \(x\in\Delta_2^\circ\) arises from a positive diagonal \(D\), uniquely up to common positive scale.
\end{lemma}

\begin{proof}
The characteristic coefficients are
\[
s=\sum_i p_id_i,
\qquad
b_D=\sum_{i<j}m_{ij}d_id_j,
\qquad
c_D=qd_1d_2d_3.
\]
Substituting \(d_i=sx_i/p_i\) gives
\[
b_D=s^2B_\beta(x),
\qquad
c_D=s^3\kappa x_1x_2x_3.
\]
Conversely \(d_i=sx_i/p_i\) reconstructs \(D\) for any \(s>0\).
\end{proof}

The radial substitution \(\lambda=sz\) shows that \(s\) changes only eigenvalue magnitudes, not arguments. Thus the full positive-diagonal orbit modulo common scale is exactly the open simplex. The four ratios \((\beta,\kappa)\) are unchanged by positive left-diagonal scaling.
